\documentclass{template}

\usepackage{mathtools}
\pgfplotsset{compat=1.18}

\newcommand{\qedwhite}{\hfill \ensuremath{\Box}}


\begin{document}

\begin{center}
    \Large \textcolor{blue!70!black}{\textbf{Commutative Algebra Problem Set \#4}} \\
    \normalsize \textit{Dmytro Zakharov}
\end{center}

\begin{problem}[Exercise 1]
    Let $A$ be a domain and let $M$ be a finitely presented $A$-module. Prove that 
    $M$ is flat (equivalently, projective or locally free) if and only if the 
    function
    \begin{equation*}
        \delta_M: \text{Spec}(A) \to \mathbb{N}, \quad \delta_M(x) = \dim_{\kappa(x)}M(x)
    \end{equation*}
    is constant. Find a counterexample to the ``if'' part with $A$ not a domain.
\end{problem}

\textbf{Solution.} $(\Rightarrow)$. Assume $M$ is flat and f.p. (and thus locally free), then localization
$M_x \cong A_x^{n_x}$ for every $x \in \text{Spec}(A)$ and some $n_x \in
\mathbb{Z}_{\geq 0}$. Let us find $\delta_M(x)$:
\begin{align*}
    \delta_M(x) &= \text{dim}_{\kappa(x)}M(x) = \text{dim}_{\kappa(x)}(M \otimes_A \kappa(x)) \\
    &= \text{dim}_{\kappa(x)}(M_x \otimes_{A_x} \kappa(x)) \\
    &= \text{dim}_{\kappa(x)}(A_x^{n_x} \otimes_{A_x} \kappa(x)) \\
    &= \text{dim}_{\kappa(x)}(\kappa(x)^{n_x}) = n_x.
\end{align*}

Thus we shall show that $M_x \cong A_x^n$ for some constant $n \in
\mathbb{Z}_{\geq 0}$ not depending on $x \in \text{Spec}(A)$. Let $K :=
\text{Frac}(A)$. We claim that $n = \text{dim}_KM_K$. Indeed, note that $M_K
\cong (M_x) \otimes_{A_x} K \cong K^{n_x}$. By taking $\text{dim}_K(\cdot)$ of
both sides, we get $n_x=n$, which completes this part.

$(\Leftarrow)$. Now assume $\delta_M(x) = n$ for all $x \in \text{Spec}(A)$. We
shall show that $M$ is locally free. As noted in the hint, replace $A$ with the
local ring $A_x$ and $M$ with $M_x$ for arbitrary $x \in \text{Spec}(A)$. Note
that the function $\delta_{M_x}$ is also trivially constant. 

Let us drop the subscript $x$ and assume that $(A,\mathfrak{m})$ is a local
domain with maximal ideal $\mathfrak{m}=s \in \text{Spec}(A)$. Thus, we have
\begin{equation*}
    \delta_M(s) = \text{dim}_{A/\mathfrak{m}}(M \otimes_A A/\mathfrak{m}) = \text{dim}_{A/\mathfrak{m}}(M/\mathfrak{m}M) = n.
\end{equation*}

Using Nakayama's lemma, we infer that $M$ is generated by $n$ elements, thus
defining the following exact sequence:
\begin{equation*}
    0 \to N \to A^n \xrightarrow{f} M \to 0, \quad N = \text{ker}(f)
\end{equation*}

Now apply $(\cdot)\otimes_A K$ functor (and it is exact because $K$ is a flat
$A$-module since it is a localization of $A$):
\begin{equation*}
    0 \to N \otimes_A K \to K^n \xrightarrow{f \otimes 1} M \otimes_A K \to 0
\end{equation*}

Now note that $\text{dim}_K(M \otimes_A K) = n$ and $f \otimes 1$ is a
surjective linear map from $K^n$ to a $K$-vector space of dimension $n$. This
implies that $f \otimes 1$ is an isomorphism and thus $N \otimes_A K = 0$. This
implies that $N=0$ and thus original map $f: A^n \to M$ is an isomorphism. Thus,
$M \cong A^n$, so $M$ is free.

Now recall that by $M$ we actually meant $M_x$, so we got $M_x \cong (A_x)^n$.
We conclude that $M$ is locally free. \qedwhite

\textbf{Counterexample.} Let us construct a ring $A$ that is not a domain,
a finitely presented $A$-module $M$, such that $\delta_M(x)$ is constant, 
but $M$ is not flat.

Take $A := k[T]/(T^2)$ and $M := k = A/(T)$. $A$ is clearly not a domain and $M$
is finitely presented. Now note that $\text{Spec}(A) = \{\mathfrak{m}\}$ for
$\mathfrak{m}=(T)$, so the function $\delta_M$ has only one input:
\begin{equation*}
    \delta_M(\mathfrak{m}) = \text{dim}_{\kappa(\mathfrak{m})}(M \otimes_A \kappa(\mathfrak{m})) = \text{dim}_k(k \otimes_A k) = 1.
\end{equation*}

Let us show that $M$ is not flat. We need to find some ideal $I$ such that $I
\hookrightarrow A$ is injective but $\sigma: M \otimes_A I \to M \otimes_A A$ is
not. Indeed, simply take $I := (T)$. Take non-zero element $z := 1 \otimes T \in
M \otimes_A I$. $\sigma$ sends this element to $1 \otimes T$, but in $M
\otimes_A A$. But inside $M \otimes_A A$, $1 \otimes T = T \otimes 1$, which 
is $0$.

\begin{problem}[Exercise 2]
    Prove that the following ring homomorphisms are not flat.
    \begin{enumerate}
        \item[(a)] $k[x^2,x^3] \hookrightarrow k[x]$;
        \item[(b)] $k[x^2,xy,y^2] \hookrightarrow k[x,y]$;
        \item[(c)] $k[x,xy] \hookrightarrow k[x,y]$.  
    \end{enumerate}
\end{problem}

\textbf{Solution.} For all parts for $A \hookrightarrow B$ we shall show that the
functor $(\cdot) \otimes_A B$ is not left exact, that is for some ideal $I \subseteq A$
\begin{equation*}
    \{ 0 \to I \to A \; \text{exact} \} \; \text{such that} \; \{0 \to I \otimes_A B \to B\} \; \text{is not exact},
\end{equation*}

where the map $\sigma: I \otimes_A B \to B$ is given by $i \otimes b \mapsto ib$.

\textit{Part (a).} Take $I := (x^2,x^3)$. We claim that
\begin{equation*}
    z := x^2 \otimes x - x^3 \otimes 1 \in \text{ker}(\sigma) \setminus \{0\}.
\end{equation*}

Clearly $z \in \text{ker}(\sigma)$ since $\sigma(z) = x^3 - x^3 = 0$. However, I
honestly cannot rigorously show that $z$ is non-zero. Intuition is the
following: $z$ were zero if we could ``factor out'' $x$ somehow. However, $x$
\textit{is not an element of $k[x^2,x^3]$}, thus this operation is impossible.
There is no other way how we can get $0$ and thus this element is non-zero.

\textit{Part (b).} Take $I := (x^2,xy)$. Then, 
\begin{equation*}
    z := x^2 \otimes y - xy \otimes x \in \text{ker}(\sigma) \setminus \{0\}.
\end{equation*}

\textit{Part (c).} Take $I := (x,xy)$. Then, 
\begin{equation*}
    z := x \otimes y - xy \otimes 1 \in \text{ker}(\sigma) \setminus \{0\}.
\end{equation*}

\begin{problem}[Exercise 3]
    Let $A$ be an $\mathbb{F}_p$-algebra and define the \textbf{Frobenius
    morphism} of $A$ as $\pi: A \to A, x \mapsto x^p$. Prove that $\pi$ induces
    the identity map $\text{Spec}(A) \to \text{Spec}(A)$.
\end{problem}

\textbf{Proof.} By definition, a ring homomorphism $f: A \to B$ induces a
continuous map $f^*: \text{Spec}(B) \to \text{Spec}(A)$ defined by
$f^*(\mathfrak{q}) = f^{-1}(\mathfrak{q})$ for any $\mathfrak{q} \in
\text{Spec}(B)$. In our case, we have a ring homomorphism $\pi: x \mapsto x^p$,
and we shall show that $\pi^*(\mathfrak{p}) = \mathfrak{p}$ for all
$\mathfrak{p} \in \text{Spec}(A)$.

Note that by definition $\pi^*(\mathfrak{p}) = \pi^{-1}(\mathfrak{p}) = \{a \in
A: a^p \in \mathfrak{p}\}$. But note that $\mathfrak{p}$ is prime, so $a^p \in
\mathfrak{p}$ automatically implies $a \in \mathfrak{p}$. The reverse
implication is trivial. Thus, $\pi^*(\mathfrak{p}) = \mathfrak{p}$. \qedwhite


\begin{problem}[Exercise 4]
    This problem consists of two parts:
    \begin{enumerate}
        \item[(a)] Let $A=\mathbb{F}_p[T_1,\dots,T_n]$. Prove that $\pi$ is finite and flat.
        \item[(b)] Let $A=\mathbb{F}_p[T^2,T^3]$. Prove that $\pi$ is finite but not flat. 
    \end{enumerate}
\end{problem}

\textbf{Solution.} 

\textit{Part (a).} Let us show that $A=\mathbb{F}_p[T_1,\dots,T_n]$ is a finite
and flat $\mathbb{F}_p[T_1^p,\dots,T_n^p]$-module. Notice that it is finite.
Indeed, consider arbitrary monomial $T_1^{\alpha_1}\dots T_n^{\alpha_n}$. Divide
each $\alpha_i$ by $p$: $\alpha_i = p\beta_i + \gamma_i$ for $0 \leq \gamma_i <
p$. Thus:
\begin{equation*}
    T_1^{\alpha_1}\ldots T_n^{\alpha_n} = (T_1^{\beta_1}\ldots T_n^{\beta_n})^p \cdot (T_1^{\gamma_1}\ldots T_n^{\gamma_n})
\end{equation*}

Note that $(T_1^{\beta_1}\ldots T_n^{\beta_n})^p \in
\mathbb{F}_p[T_1^p,\ldots,T_n^p]$, and the remaining term is in
\begin{equation*}
    S := \{T_1^{\gamma_1}\ldots T_n^{\gamma_n}: 0 \leq \gamma_i < p \; \text{for all $i$}\}
\end{equation*}

Thus, $S$ generates $A$ as an $A'$-module. Since $\#S = p^n$, $\pi$ is finite.
Moreover, $A$ is a free $A'$-module since $S$ forms a basis (trivially, all
elements of $S$ are linearly independent over
$\mathbb{F}_p[T_1^p,\dots,T_n^p]$), so in particular $\pi$ is flat.

\textit{Part (b).} The proof of $\pi$ being finite is the same as for 
\textit{part (a)}, where the set $S$ is given by:
\begin{equation*}
    S = \{T^{2i}T^{3j}: 0 \leq i,j < p\}.
\end{equation*}

However, in this case $S$ does not form a basis. As in Exercise 2, we shall find
ideal $I \subseteq \mathbb{F}_p[T^{2p},T^{3p}]$ such that $\sigma: I \otimes
\mathbb{F}_p[T^2,T^3] \to \mathbb{F}_p[T^2,T^3]$ given by $a(T^p) \otimes b(T)
\mapsto a(T^p)b(T)$ is not injective. Indeed, take $I=(T^{2p},T^{3p})$. Notice
that $z := T^{2p} \otimes (-T^p) + T^{3p} \otimes 1 \in \ker(\sigma)
\setminus \{0\}$. \qedwhite

\begin{problem}[Exercise 5]
    Let $A=k[x,y,z]/(x^2+y^2+z^2-1)$ where $k \not\supseteq \mathbb{F}_2$. Consider the 
    module of Kähler differentials $M=\Omega^1_{A/k}$. Show that $M$ is locally 
    free (or equivalently projective) and describe it as a direct summand of $A^3$.
\end{problem}

\textbf{Solution.} Recall that for $S=k[x_1,\dots,x_n]$, ideal $I=(f_1,\dots,f_r)
\subseteq S$, and $A = S/I$, we have
\begin{equation*}
    \Omega_{A/k} \cong \frac{\Omega_{S/k}}{I\Omega_{S/k} + S(d(i): i \in I)}.
\end{equation*}

We have $S=k[x,y,z]$ and $I=(\Gamma)$ for $\Gamma(x,y,z)=x^2+y^2+z^2-1$. Thus:
\begin{align*}
    \Omega_{A/k} &\cong \frac{\Omega_{S}}{\Gamma\Omega_{S} + Sd\Gamma} \cong \frac{\Omega_S / \Gamma\Omega_S}{(\Gamma\Omega_S+Sd\Gamma)/(\Gamma\Omega_S)} \cong \frac{A^3}{Ad\Gamma},
\end{align*}

where $d\Gamma = 2xdx+2ydy+2zdz$. So let us develop some intuition of what we
got so far. Note that element is zero in $\Omega_{A/k}$ iff it is of form
$f(2xdx+2ydy+2zdz)$, which intuitively are vectors, normal to the sphere at
$(x,y,z)$ (scaled by some function $f \in A$). Thus, we expect $\Omega_{A/k}
\cong A^3/d\Gamma$ to be the module of all tangent vector fields on the sphere.

Namely, suppose $\mathbf{u} = (u_x,u_y,u_z) \in A^3$ be any vector field and let
$\boldsymbol{\nu}=(x,y,z)$ be the normal vector. Define the natural scalar
product of $\mathbf{u}$ and $\boldsymbol{\nu}$ as $s: A^3 \to A$. Then, the
projection of $\mathbf{u}$ onto $\boldsymbol{\nu}$ is $\pi: A^3 \to A^3$,
$\mathbf{u} \mapsto s(\mathbf{u})\boldsymbol{\nu}$. Similarly, define the
tangent component as $\tau: A^3 \to A^3$ as $\mathbf{u} \mapsto \mathbf{u} -
s(\mathbf{u})\boldsymbol{\nu}$. Thus we see:
\begin{equation*}
    \mathbf{u} = \pi(\mathbf{u}) + \tau(\mathbf{u})
\end{equation*}

Now, suppose $N = \text{im}(\pi)$ and $T := \text{im}(\tau)$. Via the
construction above, we have shown that $A^3 = N + T$. But clearly, $N \cap T =
\{0\}$, so we have $A^3 = N \oplus T$. But notice that $N = \{f\boldsymbol{\nu}:
f \in A\} \cong A$ and $T = \text{im}(\tau) = \text{ker}(\pi)$. Thus:
\begin{equation*}
    \Omega_{A/k} \cong \frac{A^3}{\text{im}(\pi)} \cong \ker(\pi) \; \text{and thus} \; \boxed{A^3 = A \oplus \Omega_{A/k}}
\end{equation*}

This shows that $\Omega_{A/k}$ is projective.\qedwhite

\end{document}